\chapter{Firmware}

Il file principale \texttt{Pinacoteca.ino} istanzia i moduli e definisce i pin. La logica e organizzata in due fasi ben distinte, in modo da avere un ciclo di esecuzione prevedibile e facilmente verificabile.

\section{Setup}
Durante \texttt{setup()} vengono inizializzati:
\begin{itemize}
  \item seriale per debug e diagnostica;
  \item servo e tornello con posizione iniziale nota;
  \item semaforo, termostato, illuminazione e umidita;
  \item pannello LCD con messaggio di avvio.
\end{itemize}
La scelta di inizializzare prima il servo e poi gli altri moduli riduce il rischio di movimenti non voluti durante la configurazione dei pin.

\section{Loop}
Nel ciclo principale, ogni modulo aggiorna il proprio stato in una sequenza coerente:
\begin{enumerate}
  \item aggiornamento tornello e conteggio persone;
  \item aggiornamento semaforo in base al limite;
  \item regolazione clima (riscaldamento o raffreddamento);
  \item regolazione illuminazione;
  \item controllo umidita;
  \item refresh display con dati correnti.
\end{enumerate}
L'ordine e intenzionale: si acquisiscono prima gli input umani, poi si aggiornano gli indicatori di stato e infine si aggiorna l'interfaccia LCD. In questo modo il display riflette l'ultima decisione del ciclo.

\section{Gestione errori}
I moduli che leggono sensori controllano la validita della misura. Se l'input non e valido, le uscite sono poste in stato sicuro e la funzione di aggiornamento ritorna \texttt{false}. Questo comportamento e coerente con la filosofia fail-safe del progetto.

\section{Configurazione dei target}
I target sono passati ai costruttori dei moduli:
\begin{itemize}
  \item \texttt{Thermostat}: temperatura desiderata in gradi Celsius.
  \item \texttt{LightingControl}: illuminamento target in lux.
  \item \texttt{HumidifierControl}: umidita target in percento.
  \item \texttt{Turnstile}: numero massimo di persone.
\end{itemize}
La configurazione centralizzata nel file principale rende immediata la taratura dell'impianto senza modificare i singoli moduli.
